\documentclass[11pt]{article}
\usepackage[T1]{fontenc}
\usepackage{lmodern}
\usepackage{amsmath,amssymb}
\usepackage{booktabs}
\usepackage{array}
\usepackage[a4paper,margin=2.2cm]{geometry}
\usepackage[hidelinks]{hyperref}
\newcommand{\doi}[1]{\href{https://doi.org/#1}{\texttt{#1}}}
\newcommand{\dd}{\mathrm{d}}
\setlength{\parskip}{2pt}
\title{\vspace{-1.4cm}\textbf{Model v4 --- field equations for the three papers}\\[2pt]
\large Gate document for approval, before any package is built}
\author{Dr.\ Vipin \\ Gurugram University}
\date{30 October 2026}
\begin{document}
\maketitle
\vspace{-7mm}
\small
\noindent\textbf{Scope of this document.} The coupled field equations of the three papers, counted, with a
source for every term. The model stops before the boundary conditions, the dispersion determinant and any
solution step: those are the scholars' work. Status marks in the source columns are
[R]~read in the source, [T]~standard textbook law (no equation number needed),
[D]~form fixed by your instruction, [P]~proposed extension with no source; for every [P] the scholar
derives it from the stated textbook law. Where the entry does not quote an equation number, I did not read
one and the scholar is asked to locate it and write it down.

\section*{1. Paper A (Sumit) --- piezoelectric semiconductor plate, ZnO type}

\textbf{Unknowns (5):} $u_x$, $u_z$, $\varphi$, $\Delta n$, $\theta$.\quad
\textbf{Equations (5):} momentum $x$, momentum $z$, Gauss, carrier continuity, MGT heat.

\noindent\textbf{Geometry.} Plane strain in $x$--$z$, poling along $z$, faces $z=\pm h$, fields
proportional to $e^{i(kx-\omega t)}$.

\noindent\textbf{Strain and field variables.} $\varepsilon_{ij}=\tfrac12(u_{i,j}+u_{j,i})$ ($i,j=x,z$),
$E_i=-\varphi_{,i}$, carrier density $n=n_0+\Delta n$, temperature rise $\theta=T-T_0$.

\noindent\textbf{Constitutive relations with the size effect.} The total stress carries the nonlocal
strain gradient operator, the classical (local) stress does not:
\begin{align}
t_{ij}&=\bigl(1-l_{\rm nl}^2\nabla^2\bigr)\sigma_{ij}, \qquad
\sigma_{ij}=c_{ijkl}\bigl(1-l_g^2\nabla^2\bigr)\varepsilon_{kl}-e_{kij}E_k-\beta_{ij}\theta, \tag{A1}\\
D_i&=e_{ikl}\bigl(1-l_g^2\nabla^2\bigr)\varepsilon_{kl}+\epsilon_{ik}E_k+p_i\theta, \tag{A2}
\end{align}
with $\beta_{ij}=c_{ijkl}\alpha_{kl}$ and $p_i$ the pyroelectric vector. Only the \emph{elastic} part of
$\sigma$ carries the gradient operator; the piezoelectric, thermal and dielectric terms are not strain
gradients of the displacement and are left local. $l_{\rm nl}$ and $l_g$ enter only as $l/h$.
\begin{center}\footnotesize
\begin{tabular}{p{2.5cm}p{5.3cm}p{6.3cm}}
\toprule
term & meaning & source (status)\\
\midrule
(A1) NSGT form & nonlocal (softening) and gradient (stiffening) operators acting on the elastic stress &
Lim, Zhang \& Reddy, \emph{J.~Mech.~Phys.~Solids} \textbf{78} (2015) 298, \doi{10.1016/j.jmps.2015.02.001} [T]; no eq.\ no.\ read\\
$\beta_{ij}\theta$ & thermal stress by Duhamel's analogy & textbook thermoelasticity [T]\\
$p_i\theta$ & pyroelectric polarisation & Bursian \& Zaikovskii line of work; no eq.\ no.\ read [T]\\
\bottomrule
\end{tabular}
\end{center}

\noindent\textbf{Field equations.}
\begin{align}
t_{xj,j}&=\rho\,\ddot u_x, \qquad t_{zj,j}=\rho\,\ddot u_z, \tag{A3}\\
D_{i,i}&=-q\,\Delta n, \tag{A4}\\
\frac{\partial \Delta n}{\partial t}&=D_n\nabla^2\Delta n+\mu_n n_0\,\nabla\!\cdot\!E
+v_d\,\frac{\partial \Delta n}{\partial x}+\frac{\mu_n n_0}{q}\,\beta_{\rm th}\nabla^2\theta, \tag{A5}\\
\Bigl(1+\tau_q\frac{\partial}{\partial t}\Bigr)
\Bigl[\rho c_p\dot\theta+T_0\beta_{ij}\dot\varepsilon_{ij}
+T_0 p_i \dot E_i-\frac{E_{g}}{q}\,\dot{(\Delta n)}\Bigr]&=k_{\rm th}\nabla^2\theta . \tag{A6}
\end{align}
\begin{center}\footnotesize
\begin{tabular}{p{2.5cm}p{5.4cm}p{6.2cm}}
\toprule
term & meaning & source (status)\\
\midrule
(A3) & balance of momentum, total stress & textbook [T]\\
(A4) & Gauss law with the mobile carrier charge & textbook; PSC form in Yang \& Zhou, \emph{Int.\ J.\ Solids Struct.} \textbf{42} (2005) 3171, \doi{10.1016/j.ijsolstr.2004.10.011} [T]\\
(A5) drift term $v_d\partial_x\Delta n$ & drift current from a bias field $E_0$, $v_d=\mu_nE_0$ & Yang \& Zhou (2005), as above [T]; the tool is published, the combination below is not\\
(A5) thermodiffusion term & carrier flux driven by a temperature gradient (Seebeck), $\beta_{\rm th}$ the thermoelectric coefficient & textbook irreversible thermodynamics [T]\\
(A6) MGT operator & \emph{(base form, $K^*=0$, one time derivative cancelled)} $(1+\tau_q\partial_t)$ acting on the whole bracket & form fixed by your instruction [D]; attendant paper: Gupta et al., \emph{Int.\ J.\ Numer.\ Methods Heat Fluid Flow} \textbf{35}(9) (2025) 3400, \doi{10.1108/HFF-04-2025-0254}\\
(A6) $T_0\beta_{ij}\dot\varepsilon_{ij}$ & thermoelastic (dilatational) coupling & textbook [T]\\
(A6) $T_0p_i\dot E_i$ & electrocaloric coupling, the reverse of pyroelectricity & textbook [T]\\
(A6) $-(E_g/q)\dot{\Delta n}$ & carrier recombination/energy release as a heat source, the direct carrier--temperature coupling & [P]: derived from the energy balance of the electron gas; closest published PSC wave paper (Yang \& Zhou 2005) keeps the mechanical and electrical fields uncoupled from heat\\
\bottomrule
\end{tabular}
\end{center}

\noindent\textbf{What makes this paper different from the closest papers.} The carrier--temperature loop is
closed in both directions: the temperature gradient drives a carrier current through (A5), and the carrier
perturbation heats the lattice through (A6). Published PSC plate papers close at most one of the two. The
risk is that the loop is small: both coefficients are proportional to $T_0$ times a transport coefficient,
so the effect may sit below the root-finding resolution. How the paper handles that is in section 4.

\section*{2. Paper B (Neha) --- 1D hexagonal piezoelectric quasicrystal plate}

\textbf{Unknowns (5):} $u_x$, $u_z$, $w_z$, $\varphi$, $\theta$, with $w_z$ the phason displacement along
the quasiperiodic direction.\quad
\textbf{Equations (5):} phonon momentum $x$ and $z$, phason balance, Gauss, MGT heat.

\noindent\textbf{Constitutive relations.} Phonon stress carries the size effect; phason stress is carried
both ways, since which of the two the referee will accept is not settled:
\begin{align}
t_{ij}&=\bigl(1-l_{\rm nl}^2\nabla^2\bigr)\sigma_{ij},\qquad
\sigma_{ij}=C_{ijkl}\bigl(1-l_g^2\nabla^2\bigr)\varepsilon_{kl}
+R_{ijkl}w_{kl}-e_{kij}E_k-\beta_{ij}\theta, \tag{B1}\\
H_{ij}&=R_{klij}\varepsilon_{kl}+K_{ijkl}w_{kl}-\beta^{w}_{ij}\theta, \tag{B2}\\
D_i&=e_{ikl}\bigl(1-l_g^2\nabla^2\bigr)\varepsilon_{kl}+\epsilon_{ik}E_k+p_i\theta, \tag{B3}
\end{align}
with $w_{kl}=\partial w_k/\partial x_l$ and $\beta^w$ the thermo-phason coupling. The phason stress in (B2)
is kept classical in the main model; the NSGT-on-phason variant is a labelled sensitivity result only,
because no published work applies the operator to the phason strain and there is no measured phason
correlation length.
\begin{center}\footnotesize
\begin{tabular}{p{2.4cm}p{5.4cm}p{6.3cm}}
\toprule
term & meaning & source (status)\\
\midrule
$R_{ijkl}$, $K_{ijkl}$ & phonon--phason coupling and phason elastic constants, 1D hexagonal QC & benchmark material set of \emph{Acta Mech.} \textbf{232} (2021) 135--151, \doi{10.1007/s00707-020-02811-7} [T]; values to be transcribed by the scholar\\
$\beta^{w}_{ij}\theta$ & thermo-phason coupling & [P] extension of Duhamel's analogy to the phason stress; no source found\\
(B1) NSGT & as in (A1) & Lim, Zhang \& Reddy (2015), as above [T]\\
\bottomrule
\end{tabular}
\end{center}

\noindent\textbf{Field equations.}
\begin{align}
t_{xj,j}&=\rho\,\ddot u_x,\qquad t_{zj,j}=\rho\,\ddot u_z, \tag{B4}\\
H_{zj,j}&=\rho_p\,\ddot w_z+\gamma\,\dot w_z, \tag{B5}\\
D_{i,i}&=0, \tag{B6}\\
\Bigl(1+\tau_q\frac{\partial}{\partial t}\Bigr)
\Bigl[\rho c_p\dot\theta+T_0\beta_{ij}\dot\varepsilon_{ij}
+T_0\beta^{w}_{ij}\dot w_{ij}+T_0 p_i\dot E_i\Bigr]&=k_{\rm th}\nabla^2\theta . \tag{B7}
\end{align}
\begin{center}\footnotesize
\begin{tabular}{p{2.4cm}p{5.4cm}p{6.3cm}}
\toprule
term & meaning & source (status)\\
\midrule
(B5) inertia $\rho_p\ddot w_z$ & wave-like phason motion & Lubensky, \emph{Phys.\ Rev.\ B} \textbf{32} (1985) 7444, \doi{10.1103/PhysRevB.32.7444} [T]\\
(B5) relaxation $\gamma\dot w_z$ & diffusive (overdamped) phason motion; the two limits are one equation with $\gamma\to0$ and $\rho_p\to0$ & same source [T]; the comparison of the two limits in a plate is not published\\
(B6) & Gauss law, no free carriers in the dielectric & textbook [T]\\
(B7) MGT operator & as (A6) & [D]\\
(B7) $T_0\beta^w_{ij}\dot w_{ij}$ & work of the phason stress against the temperature & [P]; the heat equation in QC papers is written without it\\
\bottomrule
\end{tabular}
\end{center}

\noindent\textbf{What makes this paper different.} One phason equation carries both inertia and relaxation,
so the wave-like and the diffusive descriptions become the two ends of a single parameter instead of two
competing models, and the heat equation is coupled to the phason work as well as the phonon work. The risk
is the opposite of Paper A's: the relaxation term is easy to see in the numbers but hard to justify from a
measurement, so $\gamma$ must be swept and reported, not fixed.

\section*{3. Paper C (Anu) --- magneto-electro-elastic plate, BaTiO$_3$/CoFe$_2$O$_4$}

\textbf{Unknowns (5):} $u_x$, $u_z$, $\varphi$, $\psi$, $\theta$, with $\psi$ the magnetic potential,
$H_i=-\psi_{,i}$.\quad
\textbf{Equations (5):} momentum $x$, momentum $z$, electric Gauss, magnetic Gauss, MGT heat.

\noindent\textbf{Constitutive relations.}
\begin{align}
t_{ij}&=\bigl(1-l_{\rm nl}^2\nabla^2\bigr)\sigma_{ij},\qquad
\sigma_{ij}=C_{ijkl}\bigl(1-l_g^2\nabla^2\bigr)\varepsilon_{kl}-e_{kij}E_k-q_{kij}H_k-\beta_{ij}\theta, \tag{C1}\\
D_i&=e_{ikl}\bigl(1-l_g^2\nabla^2\bigr)\varepsilon_{kl}+\epsilon_{ik}E_k+\alpha_{ik}H_k+p_i\theta, \tag{C2}\\
B_i&=q_{ikl}\bigl(1-l_g^2\nabla^2\bigr)\varepsilon_{kl}+\alpha_{ik}E_k+\mu_{ik}H_k+m_i\theta, \tag{C3}
\end{align}
with $q_{kij}$ piezomagnetic, $\alpha_{ik}$ the magnetoelectric tensor, $p_i$ pyroelectric and $m_i$
pyromagnetic.
\begin{center}\footnotesize
\begin{tabular}{p{2.4cm}p{5.4cm}p{6.3cm}}
\toprule
term & meaning & source (status)\\
\midrule
(C1)--(C3) linear MEE constitutive law & full coupled set for a piezoelectric--piezomagnetic composite & Chen, Pan \& Chen, \emph{Int.\ J.\ Solids Struct.} \textbf{44} (2007) 1073, \doi{10.1016/j.ijsolstr.2006.06.003} [T]; no eq.\ no.\ read\\
$p_i\theta$, $m_i\theta$ & pyroelectric and pyromagnetic terms & textbook extension [T]; rarely retained in plate papers, which is the opening\\
$\alpha_{ik}$ direct ME coefficient & product property of the two phases; in a laminate it is not a material constant & reported as an effective quantity from homogenisation [T], not measured & \\
\bottomrule
\end{tabular}
\end{center}

\noindent\textbf{Field equations.}
\begin{align}
t_{xj,j}&=\rho\ddot u_x,\qquad t_{zj,j}=\rho\ddot u_z, \tag{C4}\\
D_{i,i}&=0,\qquad B_{i,i}=0, \tag{C5}\\
\Bigl(1+\tau_q\frac{\partial}{\partial t}\Bigr)
\Bigl[\rho c_p\dot\theta+T_0\beta_{ij}\dot\varepsilon_{ij}
+T_0p_i\dot E_i+T_0m_i\dot H_i\Bigr]&=k_{\rm th}\nabla^2\theta . \tag{C6}
\end{align}
\begin{center}\footnotesize
\begin{tabular}{p{2.4cm}p{5.4cm}p{6.3cm}}
\toprule
term & meaning & source (status)\\
\midrule
(C5) & electric and magnetic Gauss laws in a dielectric with no free charge or current & textbook [T]\\
(C6) $T_0m_i\dot H_i$ & magnetocaloric coupling, the reverse of pyromagneticity & textbook [T]\\
homogenisation & effective constants of the laminate from the two phases; two schemes carried & Mori--Tanaka and rule-of-mixtures, textbook micromechanics [T]\\
\bottomrule
\end{tabular}
\end{center}

\noindent\textbf{What makes this paper different.} The thermal field is coupled to all three of the
electric, magnetic and mechanical fields at once, and the magnetoelectric coefficient is treated as an
effective, frequency-dependent quantity obtained from homogenisation rather than as a material constant.
Published MEE plate papers either drop the pyromagnetic term or fix $\alpha$ at a measured value. The risk
is that the two homogenisation schemes disagree more than the effect being studied, which the paper must
show on one figure before anything else.

\section*{4. Honest note on where the earlier findings still bite}

MGT is now part of all three models, as instructed. The earlier work in this project established two
facts that do not go away, and the papers should present them as results. First, in an unbounded medium
both Fourier and MGT conduction make the thermoelastic attenuation \emph{saturate}; MGT rescales the
plateau by $\tau_T^2/(\tau_T-\tau_q)^2$ and moves the peak from $1/\tau_T$ to $1/|\tau_T-\tau_q|$, and it
does not remove the floor. Second, whether the relaxation time is \emph{visible} depends on the plate
thickness: the crossing condition is $(2h)^*=\pi\sqrt{\kappa\tau_q}$, and with the phonon mean free path
of ZnO taken from the literature the crossing sits below the thickness at which the continuum heat model
is defensible. Each paper therefore carries one figure, early, that maps the plane of thickness against
$\tau_q$ and shades the region where the thermal relaxation is resolvable, with the trust band drawn on
it. Inside the shaded region the MGT terms do the work; outside, the adiabatic and isothermal solutions
bracket the answer to a width of order the thermoelastic coupling parameter, and that bracket is reported
as the verification of the model rather than as its replacement. Written this way, a referee sees a paper
that knows where its own effect lives, which is stronger than a paper that assumes it lives everywhere.

\section*{5. Counts, for the record}

\begin{center}\small
\begin{tabular}{lcccl}
\toprule
paper & unknowns & field equations & coupled terms beyond the closest papers & extra \\
\midrule
A (Sumit) & 5 & 5 & carrier--temperature loop in both directions, pyroelectric, drift & (A1)--(A6)\\
B (Neha) & 5 & 5 & phason inertia and relaxation in one equation, thermo-phason coupling & (B1)--(B7)\\
C (Anu) & 5 & 5 & pyroelectric and pyromagnetic together, effective $\alpha$ & (C1)--(C6)\\
\bottomrule
\end{tabular}
\end{center}

\noindent\textbf{What is deliberately not in this document:} boundary conditions, interface and face
conditions, the dispersion determinant, the solution method and any numerical value. Those are the
scholars' work and the subject of the first meeting.

\end{document}
